\subsection{对称张量的乘积}

设 \(p, q \in \mathbf{N}\) ，并令 \(\mathfrak{S}_{p|q}\) 为 \(\mathfrak{S}_{p+q}\) 的由所有置换 \(\sigma \in \mathfrak{S}_{p+q}\) （它们使 N 的区间 \((1, p)\) 与 \((p+1, p+q)\) 保持稳定）组成的子群。如果 \(\sigma \in \mathfrak{S}_p\) 且 \(\sigma' \in \mathfrak{S}_q\) ，我们可以定义一个元素 \(\sigma''\) ，它属于 \(\mathfrak{S}_{p|q}\) ，通过令 \(\sigma''(n) = \sigma(n)\) 对 \(1 \leqslant n \leqslant p\) ，且 \(\sigma''(p+n) = p + \sigma'(n)\) 对 \(1 \leqslant n \leqslant q\) ；映射 \((\sigma, \sigma') \mapsto \sigma''\) 是从 \(\mathfrak{S}_p \times \mathfrak{S}_q\) 到 \(\mathfrak{S}_{p|q}\) 上的同构。

设 \(z \in \mathbf{TS}^p(\mathbf{M})\) ，\(z' \in \mathbf{TS}^q(\mathbf{M})\) ，则元素 \(z \otimes z'\) （作为 \(\mathbf{T}^{p + q}(\mathbf{M})\) 的元素）在 \(\mathfrak{S}_{p|q}\) 下不变；因此我们可以定义元素 \(\mathrm{Tr}_{\mathfrak{S}_{p + q}/\mathfrak{S}_{p|q}}(z \otimes z')\) ，它属于 \(\mathbf{TS}^{p + q}(\mathbf{M})\) 。我们将赋予 \(\mathbf{TS}(\mathbf{M})\) 一个 A-双线性乘法 \((y, y') \mapsto yy'\) ，使得对 \(p, q \in \mathbf{N}\) ，\(z \in \mathbf{TS}^p(\mathbf{M})\) ，\(z' \in \mathbf{TS}^q(\mathbf{M})\) 我们有

\[
z z ^ {\prime} = \operatorname{Tr} _ {\mathfrak {S} _ {p + q} / \mathfrak {S} _ {p | q}} (z \otimes z ^ {\prime}).\tag{4}
\]

如果 \(y \in \mathbf{TS}(\mathbf{M})\) 且 \(y' \in \mathbf{TS}(\mathbf{M})\) ，我们将称 \(yy'\) 为 \(y\) 与 \(y'\) 的对称积。族 \((\mathbf{TS}^p(\mathbf{M}))_{p \in \mathbf{N}}\) 是 \(\mathbf{N}\) 型的一个分次，即代数 \(\mathbf{TS}(\mathbf{M})\) 的分次，且 \(\mathbf{T}(\mathbf{M})\) 的单位元是 \(\mathbf{TS}(\mathbf{M})\) 的单位元。

设 \(\mathfrak{S}_{p,q}\) 为由所有满足下式的 \(\sigma \in \mathfrak{S}_{p + q}\) 组成的集合

\[
\begin{array}{l} \sigma (1) <   \sigma (2) <   \ldots <   \sigma (p) \\ \sigma (p + 1) <   \sigma (p + 2) <   \ldots <   \sigma (p + q). \end{array}
\]

映射 \((\sigma, \tau) \mapsto \sigma \tau\) （从 \(\mathfrak{S}_{p,q} \times \mathfrak{S}_{p|q}\) 到 \(\mathfrak{S}_{p+q}\) ）是双射的（I，第60页，例 2）；因此如果 \(z \in \mathbf{T}\mathbf{S}^p(\mathbf{M})\) 且 \(z' \in \mathbf{T}\mathbf{S}^q(\mathbf{M})\) ，我们有

\[
z z ^ {\prime} = \sum_ {\sigma \in \mathfrak {S} _ {p, q}} \sigma (z \otimes z ^ {\prime}).\tag{5}
\]

命题 2. --- (i) A-代数 TS(M) 是结合的、交换的且含单位元的。

\begin{enumerate}
\def\labelenumi{(\roman{enumi})}
\setcounter{enumi}{1}
\tightlist
\item
  设 \(p_1, \ldots, p_n\) 为整数 \(> 0\) ，并设 \(\mathfrak{S}_{p_1 | \ldots | p_n}\) 为所有 \(\sigma \in \mathfrak{S}_{p_1 + \ldots + p_n}\) （使下列 \(\mathbf{N}\) 的区间保持稳定）组成的集合：
\end{enumerate}

\[
\left(1, p _ {1}\right), \left(p _ {1} + 1, p _ {1} + p _ {2}\right), \dots , \left(p _ {1} + \dots + p _ {n - 1} + 1, p _ {1} + \dots + p _ {n}\right).
\]

设 \(z_{1}\in \mathbf{T}\mathbf{S}^{p_{1}}(\mathbf{M}),\ldots ,z_{n}\in \mathbf{T}\mathbf{S}^{p_{n}}(\mathbf{M})\) ，则

\[
z _ {1} z _ {2} \dots z _ {n} = \operatorname{Tr} _ {\mathfrak {S} _ {p _ {1} + \dots + p _ {n}} / \mathfrak {S} _ {p _ {1} | \dots | p _ {n}}} (z _ {1} \otimes z _ {2} \otimes \dots \otimes z _ {n}).
\]

特别地，如果 \(x_1, \ldots, x_n \in \mathbf{M}\) ，我们有 \(x_1 \ldots x_n = s(x_1 \otimes \ldots \otimes x_n)\) 。

断言(ii)对 \(n = 1\) 是显然的。假设关系式

\[
z _ {2} \dots z _ {n} = \operatorname{Tr} _ {\mathfrak {S} _ {p _ {2} + \dots + p _ {n}} / \mathfrak {S} _ {p _ {2} | \dots | p _ {n}}} (z _ {2} \otimes \dots \otimes z _ {n})
\]

已被证明，并且让我们把 \(\mathfrak{S}_{p_2 + \ldots + p_n}\) 与 \(\mathfrak{S}_{p_1 + \ldots + p_n}\) 的由所有在 \((1, p_1)\) 上限制为恒等映射的置换组成的子群等同起来。于是

\[
\begin{array}{r l} \operatorname{Tr} _ {\mathfrak {S} _ {p _ {1} | p _ {2} + \dots + p _ {n}} / \mathfrak {S} _ {p _ {1} | p _ {2} | \dots | p _ {n}}} (z _ {1} \otimes z _ {2} \otimes \dots \otimes z _ {n}) & = \\ & = z _ {1} \otimes \operatorname{Tr} _ {\mathfrak {S} _ {p _ {2} + \dots + p _ {n}} / \mathfrak {S} _ {p _ {2} | \dots | p _ {n}}} (z _ {2} \otimes \dots \otimes z _ {n}) = z _ {1} \otimes (z _ {2} \dots z _ {n}). \end{array}
\]

因此我们有

\[
\begin{array}{l} z _ {1} z _ {2} \dots z _ {n} = z _ {1} (z _ {2} \dots z _ {n}) = \\ \quad = \operatorname{Tr} _ {\mathfrak {S} _ {p _ {1} + p _ {2} + \dots + p _ {n}} / \mathfrak {S} _ {p _ {1} | p _ {2} + \dots + p _ {n}}} (z _ {1} \otimes (z _ {2} \dots z _ {n})) \\ \quad = \operatorname{Tr} _ {\mathfrak {S} _ {p _ {1} + \dots + p _ {n}} / \mathfrak {S} _ {p _ {1} | p _ {2} + \dots + p _ {n}}} (\operatorname{Tr} _ {\mathfrak {S} _ {p _ {1} | p _ {2} + \dots + p _ {n}} / \mathfrak {S} _ {p _ {1} | p _ {2} | \dots | p _ {n}}} (z _ {1} \otimes z _ {2} \otimes \dots \otimes z _ {n})) \\ \quad = \operatorname{Tr} _ {\mathfrak {S} _ {p _ {1} + \dots + p _ {n}} / \mathfrak {S} _ {p _ {1} | \dots | p _ {n}}} (z _ {1} \otimes z _ {2} \otimes \dots \otimes z _ {n}) \end{array}
\]

这由 IV 第42页命题 1 (ii) 得出。因此 (ii) 得证。

特别地，

\[
z _ {1} \left(z _ {2} z _ {3}\right) = \operatorname{Tr} _ {\mathfrak {S} _ {p _ {1} + p _ {2} + p _ {3}} / \mathfrak {S} _ {p _ {1} | p _ {2} | p _ {3}}} \left(z _ {1} \otimes z _ {2} \otimes z _ {3}\right),
\]

并且以同样的方式我们证明

\[
\left(z _ {1} z _ {2}\right) z _ {3} = \operatorname{Tr} _ {\mathfrak {S} _ {p _ {1} + p _ {2} + p _ {3}} / \mathfrak {S} _ {p _ {1} | p _ {2} | p _ {3}}} \left(z _ {1} \otimes z _ {2} \otimes z _ {3}\right).
\]

因此代数 TS(M) 是结合的。

设 \(\sigma\) 为 \(\mathfrak{S}_{p_1 + p_2}\) 中使得

\[
\begin{array}{l} \sigma (1) = p _ {2} + 1, \sigma (2) = p _ {2} + 2, \dots , \sigma (p _ {1}) = p _ {2} + p _ {1}, \\ \sigma (p _ {1} + 1) = 1, \sigma (p _ {1} + 2) = 2, \dots , \sigma (p _ {1} + p _ {2}) = p _ {2}. \end{array}
\]

则

\[
\begin{array}{l} z _ {2} z _ {1} = \operatorname{Tr} _ {\mathfrak {S} _ {p _ {1} + p _ {2}} / \mathfrak {S} _ {p _ {2} | p _ {1}}} (z _ {2} \otimes z _ {1}) \\ = \operatorname{Tr} _ {\mathfrak {S} _ {p _ {1} + p _ {2}} / \sigma \mathfrak {S} _ {p _ {1} | p _ {2}} \sigma^ {- 1}} \sigma (z _ {1} \otimes z _ {2}) \\ = \operatorname{Tr} _ {\mathfrak {S} _ {p _ {1} + p _ {2}} / \mathfrak {S} _ {p _ {1} | p _ {2}}} (z _ {1} \otimes z _ {2}) \quad \text { by   Prop   1,(i) } \\ = z _ {1} z _ {2}. \end{array}
\]

因此代数 \(\mathbf{TS}(\mathbf{M})\) 是交换的。

应当指出，TS(M) 到 T(M) 的典范嵌入一般不是代数同态。更糟的是，TS(M) 在 T(M) 的乘法下一般不是稳定的。

